% =============================================================================
%  不对中故障的转矩与受力：MathType 可直接读取的 LaTeX 片段
%  版本：2026-11-13
%
%  用法（MathType）：
%    1) MathType → Preferences → Workspace Preferences
%       勾选 "Allow TeX language input"（允许 TeX 语言输入）；
%    2) 复制下面任意一行的公式（不含行首的 % 注释行）；
%    3) 在 MathType 编辑窗口中粘贴并回车，即转成公式。
%
%  约定：每条公式单行、独立、可单独粘贴；不使用 \begin/\end、\label、\tag、
%        \text、\operatorname 与自定义宏。
%  符号约定：体固系取 e_x=(cosβ,0,-sinβ)、e_y=(0,1,0)、e_z=(sinβ,0,cosβ)，
%            故 \hat Z^(r)=(-sinβ,0,cosβ)；若改取 β→-β，横向分量整体反号。
% =============================================================================

% A1 角不对中定义
\beta=\angle\left(z,Z\right)

% A2 实际轴在体固系中的坐标
\hat{Z}^{(r)}=\left(-\sin\beta,\ 0,\ \cos\beta\right)

% A3 姿态基（右手系）
e_x=\left(\cos\beta,\ 0,\ -\sin\beta\right),\quad e_y=\left(0,\ 1,\ 0\right),\quad e_z=\left(\sin\beta,\ 0,\ \cos\beta\right)

% A4 质心位置
r_C^{(r)}=\left(0,\ 0,\ l\right)

% A5 质心到实际轴距离
d=l\sin\beta

% A6 转角
\theta=\omega t

% B1 体固系角速度
\vec{\omega}^{(r)}=\omega\hat{Z}^{(r)}

% B2 角速度分量
\vec{\omega}^{(r)}=\omega\left(-\sin\beta,\ 0,\ \cos\beta\right)

% B3 三个分量
\omega_x=-\omega\sin\beta,\quad \omega_y=0,\quad \omega_z=\omega\cos\beta

% C1 绕质心惯量
J_C=\mathrm{diag}\left(I_t,\ I_t,\ I_p\right)

% C2 平行轴定理
J_O=J_C+m\left(l^2 E-r_C r_C^{\mathrm{T}}\right)

% C3 绕 O 的惯量
J_O=\mathrm{diag}\left(I_t+ml^2,\ I_t+ml^2,\ I_p\right)

% C4 横向等效惯量
I_t'=I_t+ml^2

% C5 简写
J_O=\mathrm{diag}\left(I_t',\ I_t',\ I_p\right)

% C6 绕轴惯量定义
I_n=\hat{Z}^{\mathrm{T}}J_O\hat{Z}

% C7 绕轴惯量展开
I_n=I_t\sin^2\beta+I_p\cos^2\beta+md^2

% C8 绕轴惯量等价式
I_n=I_t'\sin^2\beta+I_p\cos^2\beta

% D1 欧拉方程
\vec{\tau}_{ext}=J_O\dot{\vec{\omega}}+\vec{\omega}\times\left(J_O\vec{\omega}\right)

% D2 角速度
\vec{\omega}=\omega\hat{Z}

% D3 角动量
J_O\vec{\omega}=\omega\left(-I_t'\sin\beta,\ 0,\ I_p\cos\beta\right)

% D4 陀螺项
\vec{\omega}\times\left(J_O\vec{\omega}\right)=\omega^2\left(I_p-I_t'\right)\sin\beta\cos\beta\ \hat{e}_y

% D5 横向分量 x
\tau_x=-I_t'\dot{\omega}\sin\beta

% D6 横向分量 y（陀螺力矩）
\tau_y=\omega^2\left(I_p-I_t'\right)\sin\beta\cos\beta

% D7 横向分量 z
\tau_z=I_p\dot{\omega}\cos\beta

% D8 绕轴分量
\tau_Z=\hat{Z}\cdot\vec{\tau}=I_n\dot{\omega}

% E1 垂足矢量定义
r_{\perp}=r_C-\left(r_C\cdot\hat{Z}\right)\hat{Z}

% E2 垂足矢量分量
r_{\perp}=\left(l\sin\beta\cos\beta,\ 0,\ l\sin^2\beta\right)

% E3 垂足模长
\left|r_{\perp}\right|=l\sin\beta=d

% E4 离心力定义
\vec{F}_c=m\omega^2 r_{\perp}

% E5 离心力分量
\vec{F}_c=m\omega^2\left(l\sin\beta\cos\beta,\ 0,\ l\sin^2\beta\right)

% E6 离心力矩（不为零）
\vec{M}_c=r_C\times\vec{F}_c=\left(0,\ m\omega^2 l^2\sin\beta\cos\beta,\ 0\right)

% E7 陀螺项分解（含离心力矩）
\vec{\omega}\times\left(J_O\vec{\omega}\right)=\omega^2\left(I_p-I_t\right)\sin\beta\cos\beta\ \hat{e}_y-\vec{M}_c

% F1 重力在体固系
\vec{F}_g^{(r)}=\left(-mg\cos\beta\sin\theta,\ -mg\cos\theta,\ -mg\sin\beta\sin\theta\right)

% F2 重力矩
\vec{M}_g=r_C\times\vec{F}_g=\left(mgl\cos\theta,\ -mgl\cos\beta\sin\theta,\ 0\right)

% F3 绕轴重力矩
M_{gZ}=\hat{Z}\cdot\vec{M}_g=-mgd\cos\theta

% G1 质心加速度
a_{CM}=-\omega^2 r_{\perp}

% G2 力平衡
R_A+R_B=-m\omega^2 r_{\perp}-F_g-F_{coup}

% G3 合力定义
F=-m\omega^2 r_{\perp}-F_g-F_{coup}

% G4 力矩平衡
r_A\times R_A+r_B\times R_B=\vec{\omega}\times J_O\vec{\omega}-M_g-M_{coup}

% G5 合力矩定义
N=\vec{\omega}\times J_O\vec{\omega}-M_g-M_{coup}

% G6 二维叉乘恒等式
\hat{Z}\times\left(\hat{Z}\times R\right)=-R

% G7 力矩平衡的二维形式
aR_A-bR_B=\hat{Z}\times N

% G8 轴承 A 反力
R_A=\frac{bF+\hat{Z}\times N}{a+b}

% G9 轴承 B 反力
R_B=\frac{aF-\hat{Z}\times N}{a+b}

% H1 摩擦转矩模型
T_f\left(\theta,\omega\right)=T_{f0}+\mu_b r_b\left(\left|R_A\right|+\left|R_B\right|\right)+c_v\omega+T_{cage}\left(\theta\right)

% I1 绕轴转矩方程
I_n\dot{\omega}=T_e-T_L\left(\theta\right)-T_f\left(\theta,\omega\right)+M_{gZ}\left(\theta\right)+T_{mis}\left(\theta\right)

% I2 开环转速纹波
\Delta\omega_k=-\frac{T_k}{I_n k\omega}\sin\left(k\theta-\psi_k\right)

% I3 闭环转速纹波
\Delta\omega_k=\left|S\left(jk\omega\right)\right|\frac{T_k}{I_n k\omega}

% J1 静不平衡离心力
\vec{F}_u=me\omega^2

% J2 陀螺力矩（标量）
\tau_{gyro}=\omega^2\left(I_t'-I_p\right)\sin\beta\cos\beta

% J3 Hooke 联轴器角速度比
\frac{\omega_2}{\omega_1}=\frac{\cos\beta}{1-\sin^2\beta\cos^2\theta}

% J4 小角度展开（2× 来源）
\frac{\omega_2}{\omega_1}\approx 1+\frac{\beta^2}{2}\cos 2\theta

% K1 开环纹波公式
\Delta\omega_1=\frac{T_1}{I_n\omega}

% K2 台架开环数值
\Delta\omega_1=\frac{0.0571}{1.89\times10^{-5}\times62.83}\approx 459\ \mathrm{rpm}

% K3 速度环衰减
\left|S\left(j2\pi\cdot 10\right)\right|\approx 0.048

% K4 型线基波
T_1\approx\frac{0.095}{1.664}=0.0571\ \mathrm{N\cdot m}

% K5 重力矩幅值
M_{gZ}=mgd

% K6 陀螺力矩量级
\tau_{gyro}\approx 6.5\times10^{-4}\ \mathrm{N\cdot m}

% K7 离心力量级
F_c\approx 0.34\ \mathrm{N}

